\label{sec:2}

\section{Related Work}
\label{sec:2.1}
Pursuit-evasion problems - where agents compete to capture or escape one another - were first formalised through a framework for reasoning about intelligent adversaries by Isaacs\cite{isaacs1965} and remain an active area of research\cite{weintraub2020}. Well-known examples include wolf and sheep, or 2d continuous, or 3d continuous.

Determining optimal strategies in real-time or close-to real time is computationally challenging as the problem grows in dimensionality - the state-action space becomes too large to represent with traditional methods\cite{bertsekas1996,powell2007}. This ‘curse of dimensionality’\cite{bellman1957} motivates a variety of approximation approaches. McGrew et al.\cite{mcgrew2010} demonstrated an approximate dynamic programming (ADP) approach to 1v1 air combat, involving a number of additional reward-shaping techniques. While effective, this approach is limited to 2D level flight at fixed velocity, does not approach the infinite state space of continuous flight\cite{bertram3d}, and did not scale to multi-agent teams. These limitations motivate a different approach for multi-agent pursuit-evasion in a continuous 3D environment.

Bertram et al.\cite{bertram2d} proposed a novel algorithm to closely approximate the solution to a class of MDPs with sparse rewards which they term FastMDP. Instead of well-known value iteration or policy iteration methods\cite{bellman1957,sutton1998}, FastMDP reframes the MDP as a graph and computes the value function from individual reward sources, combining positive and negative peaks into a single surface. The algorithm has no dependency on the size of the state space, making it suitable for continuous spaces and real-time computation. This was initially demonstrated in a 2D UAV path planning context, showing scaling with the number of obstacles.

Further works by Bertram and Wei\cite{bertram3d} expanded FastMDP to a fully continuous 3D state space, incorporating a pseudo-6DoF kinematic model of fixed-wing aircraft and driving pursuit-evasion behaviours. That work forms the basis for the simulation and decision-making framework developed in this project.

\review{NECESSARY?}
Paragraph/justification about moving forwards with the FastMDP influence in my implementation, but what were my differentiations? Primarily reward structure, performance analysis alternative behaviour extensions\
Also could describe how my work specifialcly involves agents that are both pursuing and evading at the same time, instead of other examples like 1 pursuer and 1 evader

\section{Starting Point}
\label{sec:2.2}
I had experience with building and applying Markov models in Python (in particular Markov Chains and Hidden Markov Models) from my Part IA Machine Learning and Real-World Data course, but I had not specifically implemented a Markov Decision Process or the FastMDP algorithm before.

I did not have experience with pursuit-evasion problems or explicit kinematic models, but I had written Java code to simulate physical three-dimensional systems before with a focus on linear algebra and vector mathematics rather than aerodynamics.

I was confident with the Python programming language and the tooling around it including virtual environments, the pip package manager and the pytest testing framework. I had real-world corporate experience in managing large Python projects, and was also comfortable working with the Git version control system and GitHub actions for development workflows.

The final implementations are not based on existing codebases or machine learning libraries, but I identified similar projects from related work and examined their algorithms and pseudocode to influence my work.

\section{Theory and Design}
\label{sec:2.3}
I prepared for implementation by researching and understanding two substantial bodies of theory: Markov Decision Processes and approximation solution methods, and models of 3D kinematics. 


\subsection{Markov Decision Processes}
\label{sec:2.3.1}

A Markov Decision Process (MDP) is a mathematical framework for sequential decision-making under uncertainty. An MDP is formed of a tuple $(S, A, T, R)$ of states $S$, actions $A$, transition function $T : S \times A \times S \to [0,1]$ between states, and reward function $R : S \times A \to \mathbb{R}$ encoding desired behaviour. A policy is a mapping $\pi : S \to A$ with a value function $V^\pi : S \to \mathbb{R}$ encoding the expected reward from following that policy (with some discount factor $\gamma$ per transition), and the optimal policy $\pi^*$ is the one that maximises $V^{\pi^*}$ for all $s \in S$. Solving an MDP typically involves finding an optimal policy that covers all states.

A Bellman optimality equation can exactly represent the  optimal value function\cite{bellman1957}.

\begin{equation}
V^{\pi^*}(s) = \max_{a \in A} \left[ R(s,a) + \gamma \sum_{s' \in S} T(s, a, s')V^{\pi^*}(s') \right]
\label{eq:value_bellman}
\end{equation}

Intuitively, the best value of an optimal-policy state must equal the expected return of the best action\cite{sutton1998}. This value-iteration approach is a standard method for exact solutions, but as discussed becomes intractable for large state spaces and undefined for infinite ones. A key observation for some approximation approaches is to identify a class of MDPs with \textit{sparse rewards}. In a sparse-reward MDP, most states have reward $R(s, a) = 0$, and there may be reward peaks and wells localised around certain regions of states. In this situation, the transition function $T$ and full reward function $R$ do not need explicit construction - the value function can be approximated directly from the reward sources and bypass standard iteration techniques.

\subsection{\texttt{FastMDP} Algorithm}
\label{sec:2.3.2}

The FastMDP algorithm exploits this by constructing the value function from positive peaks and negative wells in a continuous state space. A measure of distance $\delta : S \times S \to \mathbb{R}$ must be defined for the state space to interpret the effect of each peak. For a continuous 3D environment this can be Euclidean distance. 

Each positive reward places an exponentially decaying peak, and the positive value function is defined as the $\max$ over all peaks for some state $s$.

\begin{equation}
    V^+(s) = \max_{i \in R^+} \left[ r_i \cdot \gamma_i^{\delta(s, s_i)} \right]
\end{equation}

where $R^+$ are the positive peaks, $r_i$ is the reward magnitude, $\gamma_i \in (0,1)$ is the per-source decay factor, and $\delta(s, s_i)$ is the distance measure. The negative value function is defined similarly, but negative rewards (`risk wells') are first converted into \textit{standard positive form} in order to calculate the value function identically to positive rewards before negating again.

\begin{equation}
V^-(s) = -\max_{j \in R^-} \left[ |r_j| \cdot \gamma_j^{\delta(s, s_j)} \right]
\end{equation}

Summing the positive and negative value functions gives a close approximation of solving the MDP with a full reward function\cite{bertram3d,bertram2d}. \texttt{FastMDP} also requires \textit{forward projection} functionality, where a state $s$ and action $a$ are used to compute a resulting state. As the kinematic model provides deterministic transitions $s'_a=f(s,a)$, forward projection replaces the stochastic transition function $T$ entirely. An approximation of the optimal policy can be retrieved from the value function and forward projection, and applied to determine optimal actions for any particular state.

\begin{equation}
V(s) = V^+(s) + V^-(s) \approx V^{\pi^*}
\end{equation} 
\begin{equation}
\pi(s) = \arg\max_{a \in A} V(s'_a) \approx \pi^*
\end{equation} 

\subsection{Reward Functions}
\label{sec:2.3.3}
\texttt{FastMDP} has no reliance on the size of the state space, making it well-suited to the continuous 3D environment of pursuit-evasion. At each timestep, each agent formulates its own MDP with reward sources based on its opponents and their kinematic properties. It is the parameters of these reward sources (position, magnitude, decay) that drive the agent behaviour to capture enemies and avoid being captured.

One of the project success criteria was to design novel reward functions that drive this behaviour. Bertram et al. place rewards solely based on agent position - opponents have positive peaks at their current and negative wells at projected positions, attracting pursuers and repelling collisions. I wanted to expand on this by considering velocity, heading and alignment parameters to inform reward sources. Two agents at equal distances from a target would have different tactical advantages depending on their alignment towards the target, and if their opponent is relatively aligned towards them.

This motivates incorporating \textit{velocity direction} into the reward structure. In addition to position, my MDP formulation would provide each agent with access to additional state for all opponents:
\begin{enumerate}[nosep]
    \item The unit displacement vector $\hat{r}_{pe}$ from pursuer (self) to evader (enemy).
    \item The self velocity vector $v_{self}$.
    \item The enemy velocity vector $v_{enemy}$.
    \item The scalar distance $d$.
\end{enumerate}
  
Alignment measures and velocity headings can be derived from these quantities and used to inform the rewards. The specific reward functions built from these quantities and the parameter tuning process used to refine them, are described in the Implementation chapter.

\subsection{Kinematics}
\label{sec:2.3.4}

The kinematic model follows the work of previous implementations\cite{bertram3d,park2016,huynh1987}. A pseudo 6-degrees-of-freedom model (pseudo-6DoF) was applied, where an aircraft state has three positional ($x$, $y$, $z$) coordinates and three rotational (roll, pitch, yaw) coordinates. A full 6DoF model accounts for all aerodynamic forces acting on a fixed-wing aircraft, but the pseudo-6DoF model simplifies this to a smaller set of control inputs and approximated forces for use in the equations of motion.

An aircraft is modelled by a position $(x, y, z)$ with altitude $h = z$, azimuth heading angle $\psi$, attack angle $\alpha$, roll angle $\phi$, flight path angle $\gamma$, and airspeed $V$. Position is measured in metres, speed in metres/second and all angular quantities in radians. The control inputs to an aircraft are the acceleration derived from thrust $n_x$, attack angle rate of change $\dot{\alpha}$, and roll angle rate of change $\dot{\phi}$. These inputs provide a reasonable approximation of throttle and joystick controls in a highly manoeuvrable aircraft.

The equations of motion are implemented directly from Bertram and Wei\cite{bertram3d} with two additions: kinematic bounds were applied per-agent rather than per-team and a numerical guard prevents division-by-zero of velocity in the azimuth angle calculation.

\begin{equation}
    \dot{V} = g[n_x \cos\alpha - \sin\gamma]
\end{equation}
\begin{equation}
    \dot{\gamma} = \frac{g}{V}[n_f \cos\phi - \cos\gamma]
\end{equation}
\begin{equation}
    \dot{\psi} = \frac{g \cdot n_f \sin\phi}{V \cos\gamma}
\end{equation}

where $n_f = n_x \sin\alpha + L n_f $is the normal force with lift constant $L=0.5$ and $g = 9.80665$ m/s$^2$.
The bounds were determined during implementation and enforced through clipping functions. A hard deck was implemented at $z=0$, but this was enforced as a decision penalty rather than a kinematic limit.

A pseudo-6DoF model was appropriate for the project as the aim was to examine pursuit-evasion decision-behaviour with a kinematic model of reasonable fidelty. The aerodynamic force calculations of a full-6DoF that map surface deflections to drag, lift and lateral forces were not necessary for the purpose. 

\subsection{Codebase Architecture}
\label{sec:2.3.5}
I laid out a high-level architecture for the codebase that matched with the MDP and kinematic frameworks investigated above. This had a clear separation of concerns: the simulation could be run independently; the MDP constructor took only the initial parameters required to solve it; and outputs could easily be generated from simulation state at each timestep.
The initial design was object-oriented - a central  manager constructed simulation instances that maintained lists of agents. Each agent constructed a single MDP for the entire simulation run that they updated with new values at each timestep. 

\review{figure here? i think I have some really old sketches on my ipad but they're very sloppy/incomplete? i could try and clean them up and try to make them actually demonstrate some of the original architecture but they would very much be filled-in (and i dont remember a lot of the original...)}

Early on in development this changed so that agents were constructing new MDP instances at each timestep, as there was minimal performance overhead and this better represented the dataflow of the program. The Implementation chapter further describes how the design eventually moved away from pure OOP in order to take advantage of various performance optimisations that did not interface well with Python objects. 

Python had a broad ecosystem that supported all areas of the project. NumPy provided efficient numerical compute, and a variety of libraries helped with static, dynamic, and interactive visualisations. Beyond library availability, its dynamic typing and interpreted nature allowed for more rapid prototyping whilst exploiting Numba's JIT functionality to ensure performance-heavy functions could be compiled to optimised machine code.

\section{Requirements Analysis}
\label{sec:2.4}
Before beginning the project I performed a requirements analysis using the \textit{MoSCoW} system. The \textit{Must} requirements map directly onto the (non-extension) success criteria defined in the Introduction chapter.

\begin{table}[H]
\centering
\renewcommand{\arraystretch}{1.2}
\begin{tabular}{p{0.8\textwidth}@{\hspace{1cm}}l}
\textbf{Requirement} & \textbf{Priority} \\
\midrule
3D kinematic simulation of fixed-wing aircraft using a pseudo-6DoF model & \colorbox{mustgreen}{\texttt{Must}} \\
FastMDP-based decision making framework for pursuit-evasion behaviour & \colorbox{mustgreen}{\texttt{Must}} \\
Agent actions are computed within a timeframe suitable for real-time simulation & \colorbox{mustgreen}{\texttt{Must}} \\
Valid kinematic and behavioural outputs verified against defined criteria & \colorbox{mustgreen}{\texttt{Must}} \\
\midrule
Output system producing video or visualisation results of simulation runs & \colorbox{shouldyellow}{\texttt{Should}} \\
Reward functions are configurable without changes to core implementation & \colorbox{shouldyellow}{\texttt{Should}} \\
Codebase is modular and extensible to alternative problem spaces & \colorbox{shouldyellow}{\texttt{Should}} \\
\midrule
Scale to multi-agent teams beyond 1v1 & \colorbox{couldorange}{\texttt{Could}} \\
JIT compilation via \texttt{numba} for performance gains & \colorbox{couldorange}{\texttt{Could}} \\
Demonstrate framework applicability to alternative problem spaces & \colorbox{couldorange}{\texttt{Could}} \\
\midrule
Additionally accurate kinematic modelling beyond pseudo-6DoF & \colorbox{wontred}{\texttt{Won't}} \\
Real-world hardware deployment & \colorbox{wontred}{\texttt{Won't}} \\
\end{tabular}
\caption{MoSCoW requirements analysis}
\label{tab:requirements}
\end{table}

\section{Development Methodology}
\label{sec:2.5}
The project was developed on a personal MacBook Pro (2021, Apple M1 Pro, 16GB RAM) running Python \texttt{v3.13.2}. Key libraries were verified for compatibility before implementation with particular attention to \texttt{numba}, which supports ARM-based Apple Silicon\cite{numba_install} but required testing to confirm JIT compilation behaved correctly.

Git was set up for version control throughout development, with source code, assets and dissertation drafts regularly committed to a remote GitHub repository.

An agile approach to software development was adopted throughout the project. The milestones from the project proposal were used as starting points for implementation sprints, with a short written review after each 2-3 week sprint. These reviews are included in Appendix A. This approach allowed flexibility throughout the academic year as other in-term commitments and vacations affected the development schedule, which is also reflected in the commit frequency of the repository.

Unit tests were generated for individual modules using the \texttt{pytest} framework and produced alongside each implementation. A GitHub Actions workflow was configured to run the full test suite on every remote commit, with results and coverage reports saved as build artefacts. This provided a record of test outcomes throughout development to be referenced in the Evaluation chapter. However maintaining high test coverage proved to be challenging during heavy refactoring of the central implementations and limited the effectiveness of the suite. This is discussed further in the Implementation and Evaluation chapters.

Generative AI tooling was used in two clearly delineated capacities. Firstly, it was used to generate the unit tests from the existing implementations, coming up with realistic and edge-case values for the simulation and kinematic components. Secondly, it was used to produce boilerplate code for the output and visualisation modules with the Matplotlib library. All simulation, MDP, and kinematic code was written independently. No AI-generated code was included in the main algorithmic implementations, but separate chat windows were used throughout development to aid with documentation searching and reinforcing theoretical knowledge.